\documentclass[11pt]{article}
\usepackage{mathpartir}
\usepackage{pifont}
\usepackage{multicol}
\usepackage{listings}
\usepackage{pgf}
\usepackage{tikz}
\usepackage{alltt}
\usepackage{hyperref}
\usepackage{url}
\usepackage{amsmath}
\usepackage{amssymb}
\usepackage{dafny}
\usepackage{fancyvrb}
\usetikzlibrary{arrows,automata,shapes,positioning}
\tikzstyle{block} = [rectangle, draw, fill=blue!20, 
    text width=5em, text centered, rounded corners, minimum height=2em]
\tikzstyle{bt} = [rectangle, draw, fill=blue!20, 
    text width=1em, text centered, rounded corners, minimum height=2em]
\newcommand{\xmark}{\ding{55}}

\newtheorem{defn}{Definition}
\newtheorem{crit}{Criterion}
\newcommand{\true}{\mbox{\sf true}}
\newcommand{\false}{\mbox{\sf false}}

\newcommand*\circled[1]{\tikz[baseline=(char.base)]{
            \node[shape=circle,draw,inner sep=2pt] (char) {#1};}}


\newcommand{\handout}[5]{
  \noindent
  \begin{center}
  \framebox{
    \vbox{
      \hbox to 5.78in { {\bf Software Testing, Quality Assurance and Maintenance } \hfill #2 }
      \vspace{4mm}
      \hbox to 5.78in { {\Large \hfill #5  \hfill} }
      \vspace{2mm}
      \hbox to 5.78in { {\em #3 \hfill #4} }
    }
  }
  \end{center}
  \vspace*{4mm}
}

\newcommand{\lecture}[4]{\handout{#1}{#2}{#3}{#4}{Lecture #1}}
\topmargin 0pt
\advance \topmargin by -\headheight
\advance \topmargin by -\headsep
\textheight 8.9in
\oddsidemargin 0pt
\evensidemargin \oddsidemargin
\marginparwidth 0.5in
\textwidth 6.5in

\parindent 0in
\parskip 1.5ex
%\renewcommand{\baselinestretch}{1.25}

\usepackage{enumitem}

\newtheorem{prop}{Proposition}
\newtheorem{lemma}{Lemma}
\usepackage{ebproof}
\newcommand{\qedsymbol}{\rule{1ex}{1ex}}
\newcommand{\sem}[3]{\langle #1, #2 \rangle \Downarrow #3}

\lstset{ %
language=Java,
basicstyle=\ttfamily,commentstyle=\scriptsize\itshape,showstringspaces=false,breaklines=true,numbers=left}

%\usepackage{fontspec}
%\setmonofont{Cousine}[Scale=MatchLowercase]

\begin{document}

\lecture{20X --- March 23, 2026}{Winter 2026}{Patrick Lam}{version 1}

Here's some bonus content on computing greatest common divisors, from
the \textsf{stqam-1261-dafny} repo, which you can find at
\url{https://git.uwaterloo.ca/stqam-1261/dafny}.

First, consider the greatest common divisor. We will take the textual specification to be this:
\begin{lstlisting}[language=dafny]
 /*
   Define a recursive function, called gcd, that computes a GCD of two numbers.
   Use the following properties of gcd.
   (a) gcd(n, n) == n
   (b) gcd(m, n) == gcd(m-n, n) whenever m > n
   (c) gcd(m, n) == gcd(m, n-m) whenever n > m
 */
\end{lstlisting}

We can then write this specification as a (recursive) Dafny function:
\begin{lstlisting}[language=dafny]
function gcd(m: nat, n: nat): nat
  requires m > 0 && n > 0
{
  if (m == n) then m
  else if (m > n) then gcd(m - n, n)
  else gcd(m, n - m)
}
\end{lstlisting}
For our purposes, we take this to be the definition of gcd. I don't think Dafny can prove that the
declarative definition of gcd is equal to this recursive definition. Since Dafny's natural numbers
start at 0, we specifically define \textsf{gcd} to start at 1 to ensure termination: we are subtracting
numbers and we need to subtract at least 1 to $m+n$ on each recursive call if we are to make progress.

The contract of \textsf{GcdCalc} can be as follows:
\begin{lstlisting}[language=dafny]
method GcdCalc(m: nat, n: nat) returns (res: nat)
  requires m > 0 && n > 0;
  ensures res == gcd(m, n)
\end{lstlisting}
That is, we use the \textsf{gcd} function we've defined. And we mirror the precondition of the function.

The implementation is straightforward:

\begin{lstlisting}[language=dafny]
{
  var n1, m1 := n, m;
  while (n1 != m1)
  {
    if (m1 < n1) {
     n1 := n1 - m1;
    } else {
     m1 := m1 - n1;
    }
  }
  assert(m1 == n1);
  return n1;
}
\end{lstlisting}

We don't absolutely need the final assert, but it also doesn't hurt to put it there as documentation.
What we do need, though, is an invariant. Without an invariant, Dafny doesn't know what to make of the loop.
So, here, you can read along, and agree with what I'm writing, but really, you need to practice writing
some invariants to be able to actually do this. Go do some invariant writing. More than on the assignment.

Anyway, the sensible first invariant to put in is mirroring the precondition. (Indeed, Dafny complains that
it can't show the function precondition for function \textsf{gcd} with no precondition). We need both \textsf{m}
and \textsf{n} to be positive.
\begin{lstlisting}[language=dafny]
  invariant m > 0 && n > 0
\end{lstlisting}
That does help, but Dafny still can't prove termination, nor can it prove the postcondition. Of course.
We don't know anything about what the loop is actually computing at this point, just that \textsf{m} and
\textsf{n} are strictly positive.

A good guess for the additional invariant is
\begin{lstlisting}[language=dafny]
  invariant gcd(m, n) == gcd(m1, n1)
\end{lstlisting}
I don't really know how to give you intuition about this. The
invariant has to somehow \textsf{gcd} and you should have some idea
that this gcd algorithm is preserving the gcd. For you, it can be
based on something you vaguely remember from MATH 135.
\end{document}
